\section{Introduction}
A photograph of a child holding a red ball supports several descriptions.
``A child with a red ball'' describes the scene precisely.
``Someone is outdoors'' may also be correct, but matches many images.
An image search model must learn from both descriptions while distinguishing the intended match from thousands of candidates.
Which descriptions should training reward, and how strongly?

Contrastive image--text models reward associated captions relative to competing candidates \citep{radford2021clip}.
We call these rewarded pairs \emph{positives}.
A competing caption can still correctly describe the image.
Relation annotations identify supported statements that can supply additional positives.
However, promoting a caption changes both its training role and the distribution of supervision.
The original captions share their target probability with more descriptions.
The reverse loss also averages over more caption queries.

This distinction matters when adapting an existing retrieval model.
Supported descriptions can teach useful visual distinctions without identifying the exact caption assigned to an image.
Better caption discrimination therefore need not improve source-caption retrieval.
The practical objective combines both capabilities: learn useful distinctions while retaining reliable search.

We study this objective through two linked comparisons.
First, we test caption identity while matching positive counts, initial similarity groups, learning rates, and epochs.
We repeat the comparison with a nonlinear adapter and a second pretrained encoder.
Directional interventions separate image-positive expansion from reverse-anchor expansion.
Second, we separate source-target allocation from a standard preservation penalty.
Their combination tests whether allocation contributes beyond distillation alone.
Every preservation procedure uses the same development rule and an explicit retrieval tolerance.

Prior methods convert false negatives into positives and modify negative weights \citep{byun2024mafa,jiang2023regulation}.
Composition studies already examine valid descriptions and retained capabilities \citep{kamath2024hardpositive,oh2024fsc,peleg2025clic}.
CLIP-Refine combines positive supervision with teacher outputs; WiSE-FT interpolates pretrained and adapted weights \citep{yamaguchi2025refine,wortsman2022wise}.
Our contribution concerns what caption assignment and allocation add under these preservation controls.
Both new settings reproduce the source-retrieval cost.
Caption assignment improves caption-to-image retrieval, with rate-dependent evidence in the nonlinear setting.
Matched allocation and distillation improve retrieval, but no selected AD-study family meets our joint preservation criterion.
